\section{Các giao tiếp ngoại vi}
\begin{table}[H]
\centering\small
\caption{Các giao tiếp ngoại vi được sử dụng}\label{tab:interfaces}
\begin{tabularx}{\linewidth}{@{}p{2.6cm}X>{\raggedright\arraybackslash}p{4.1cm}@{}}
\toprule
\textbf{Giao tiếp} & \textbf{Nguyên lý} & \textbf{Dùng cho} \\ \midrule
GPIO & Chân vào/ra số mức 0/1, có thể bật điện trở kéo lên/xuống & 5 cảm biến IR, cảm biến lửa (DO), còi \\
PWM (LEDC) & Khối phần cứng phát xung chu kỳ cố định, độ rộng xung (duty) điều chỉnh được; CPU không phải bật/tắt chân & Góc quay 2 servo SG90 (50\,Hz, xung 0,5–2,4\,ms) \\
SPI & Bus nối tiếp đồng bộ 4 dây (SCK, MOSI, MISO, CS), một master nhiều slave, mỗi slave một chân CS; có thể dùng DMA để truyền khối lớn & Màn hình ST7789 (40\,MHz, DMA), 2 đầu đọc MFRC522 dùng chung bus \\
ADC & Chuyển điện áp tương tự sang số (ESP8266: 10 bit ở chân A0; NodeMCU có cầu phân áp cho dải 0–3,3\,V) & Cảm biến mực nước \\
1 dây (DHT) & Giao thức riêng của DHT11 trên một chân: host kéo thấp để bắt đầu, cảm biến trả 40 bit mã hóa bằng độ rộng xung & Nhiệt độ \\
UART / USB Serial-JTAG & Truyền nối tiếp bất đồng bộ; ESP32-C3 có sẵn bộ USB Serial-JTAG để nạp firmware và xem nhật ký & Nạp chương trình, nhật ký gỡ lỗi \\
\bottomrule
\end{tabularx}
\end{table}

\begin{lydo}[Tại sao dùng ngoại vi phần cứng thay vì điều khiển chân bằng phần mềm]
PWM phần cứng giữ xung servo ổn định ngay cả khi CPU bận vẽ màn hình; SPI có DMA truyền
115\,KB ảnh màn hình mà CPU không phải chép từng byte. Nếu tạo xung bằng phần mềm, chỉ một
ngắt Wi-Fi cũng đủ làm servo rung.
\end{lydo}

\section{Giao thức mạng và thanh toán}
\begin{itemize}
  \item \textbf{Wi-Fi 802.11 b/g/n (chế độ trạm)}: ba bo cùng kết nối vào một điểm phát; địa
        chỉ IP do DHCP của router cấp.
  \item \textbf{UDP}: giao thức truyền gói không kết nối, không bảo đảm đến nơi nhưng độ trễ
        thấp và không cần bắt tay. Gói \emph{broadcast} được gửi tới mọi thiết bị trong mạng
        con, nên client không cần biết IP của máy chủ. Độ tin cậy cho sự kiện quẹt thẻ được bổ
        sung bằng ACK và gửi lại (Mục~\ref{sec:giaothuc}).
  \item \textbf{HTTP}: giao thức yêu cầu–phản hồi trên TCP, dùng cho trang web giám sát và các
        API JSON.
  \item \textbf{mDNS}: phân giải tên \code{.local} bằng multicast trong mạng cục bộ, không cần
        máy chủ DNS, cho phép mở \code{http://parking.local}.
  \item \textbf{HTTPS/TLS}: HTTP được mã hóa và xác thực máy chủ bằng chứng chỉ X.509; máy chủ
        dùng để gọi API SePay.
  \item \textbf{VietQR}: chuẩn mã QR chuyển khoản liên ngân hàng của NAPAS, dựa trên chuẩn
        EMVCo: chuỗi các trường dạng \emph{mã–độ dài–giá trị} (ngân hàng, số tài khoản, số tiền,
        nội dung) kết thúc bằng mã kiểm tra CRC16 (Mục~\ref{sec:vietqr}).
\end{itemize}
